% Supplementary Materials for STARS (Paper A).
% Compile standalone:  pdflatex supplement.tex  (twice, for cross-references)
% Submitted as a separate Supplemental file; excluded from the main-text length.
\documentclass[12pt]{article}
\usepackage[margin=1in]{geometry}
\usepackage{lmodern}
\usepackage[T1]{fontenc}
\usepackage[utf8]{inputenc}
\usepackage{microtype}
\usepackage{booktabs}
\usepackage{longtable}
\usepackage{array}
\usepackage{graphicx}
\usepackage{amsmath}
\usepackage{xcolor}
\usepackage{url}
\usepackage{setspace}
\usepackage{hyperref}
\hypersetup{colorlinks=true,linkcolor=blue,citecolor=blue,urlcolor=blue}
\onehalfspacing

% Number supplementary tables/figures with an "S" prefix.
\renewcommand{\thetable}{S\arabic{table}}
\renewcommand{\thefigure}{S\arabic{figure}}

\title{Supplementary Materials\\[2pt]
\large Perceived Stress Tracks the Tinnitus Symptom More Than the Audiometric
Threshold: A Survey-Weighted KNHANES Study (STARS)}
\author{}
\date{}

\begin{document}
\maketitle
\vspace{-3em}

\noindent These Supplementary Materials collect (i) the \emph{prospective} program
components of STARS (design only; no results in the main empirical article) and
(ii) the methods-detail tables. They are referenced from the main text as
``Supplementary Methods~S1/S2'' and ``Supplementary Table~S1--S3.''

\section*{Supplementary Methods}

\subsection*{S1. Prospective external validation (companion program)}
Cross-national transportability of a \emph{fitted} KNHANES model to NHANES
(Aim~2) is a prospective component and carries no results in the main empirical
article; the NHANES analyses reported in the main text are directional support
under a distinct distress proxy, not a validation of a transported model. To keep
the plan falsifiable and auditable, the transportability protocol is prespecified
and recorded in the repository: weighted versus unweighted performance;
calibration-in-the-large and slope before and after intercept-only then
intercept$+$slope recalibration of the frozen model; common-support restriction to
the overlapping covariate region; case-mix standardization to separate genuine
transportability from prevalence/age differences; and a reduced common model
versus country-specific extended models. Discrimination and calibration are
prespecified to be reported within subgroups (sex, age band, employment, noise
exposure, unilateral/asymmetric hearing, socioeconomic status, and race/ethnicity
where collected), with subgroup calibration and error gaps as first-class results.
These analyses are developed in the companion program of work.

\subsection*{S2. Prospective clinical extension (hospital-based cohort)}
A prospective otolaryngology/audiology cohort is specified to address what public
cross-sectional data cannot: treatment timing, sudden sensorineural hearing loss
(SSNHL) recovery trajectories, tinnitus burden over time, and hearing-aid
rehabilitation. This cohort contributes deidentified audiograms and clinical
narratives for clinical external validation and provides the substrate for the AI
extraction and safety evaluations developed in the companion SAFE-EAR paper. Raw
clinical data are not redistributed; only derived, deidentified, or synthetic
artifacts are shared through the open repository, subject to institutional review
board (IRB) approval and data-governance review.

\section*{Supplementary Tables}

\begin{table}[htbp]
\centering
\caption{Selected public datasets for HEAR-WORK AI}
\label{tab:datasets}
\begin{tabular}{p{0.18\linewidth}p{0.26\linewidth}p{0.24\linewidth}p{0.22\linewidth}}
\toprule
Dataset & Role & Strengths & Main limitations \\
\midrule
KNHANES & Primary development dataset & Korean population relevance; audiometry in selected cycles; tinnitus and health variables; survey weights & Cross-sectional; stress variables may be perceived stress rather than validated occupational stress; raw data access procedures required \\
NHANES & External validation dataset & Public CDC files; audiometry and noise exposure in selected cycles; rich health covariates; geographic validation & U.S. population; cycle-specific hearing and tinnitus variables; complex survey design \\
OHHR / Zenodo audiology data & Secondary method testing & Audiology-specific measurements; useful for model benchmarking & May lack stress, occupation, or tinnitus variables \\
OpenNeuro tinnitus datasets & Optional neurophysiology substudy & Public neuroimaging or fNIRS data for tinnitus biomarkers & Not designed for occupational stress or longitudinal treatment outcomes \\
\bottomrule
\end{tabular}
\end{table}
      % Supplementary Table S1
\begin{table}[htbp]
\centering
\caption{Target harmonized variables and their availability in the primary public datasets. Availability is cycle-dependent; each derived variable records its source items and transformation in the version-controlled data dictionary.}
\label{tab:variables}
\small
\begin{tabular}{p{0.30\linewidth}p{0.14\linewidth}p{0.22\linewidth}p{0.22\linewidth}}
\toprule
Harmonized variable & Role & KNHANES & NHANES \\
\midrule
Pure-tone thresholds (per ear, per frequency) & Outcome & Selected cycles & Selected cycles \\
Tinnitus status / bother & Outcome & Selected cycles & Selected cycles \\
Self-reported hearing difficulty & Outcome (2\textsuperscript{nd}) & Yes & Yes \\
Hearing-aid use & Outcome (2\textsuperscript{nd}) & Selected cycles & Selected cycles \\
Perceived stress / depressed mood & Exposure & Yes & Yes \\
Occupation / employment status & Modifier & Yes & Yes \\
Occupational + recreational noise & Covariate & Selected cycles & Yes \\
Sleep duration & Covariate & Selected cycles & Selected cycles \\
Smoking, alcohol & Covariate & Yes & Yes \\
Cardiometabolic disease & Covariate & Yes & Yes \\
Age, sex, education & Covariate & Yes & Yes \\
Survey weight, stratum, PSU & Design & Yes & Yes \\
\bottomrule
\end{tabular}
\end{table}
     % Supplementary Table S2
\begin{table}[htbp]
\centering
\caption{Cross-national harmonization and measurement-comparability plan (excerpt). Exact cycle-specific variable codes and response categories are pinned in the repository data dictionary. Comparability governs use: \emph{low}-comparability variables are excluded from the common transportable model and entered only in country-specific extended models. ``Sel.\ cycles'' = present in selected cycles only.}
\label{tab:harmonization}
\scriptsize
\begin{tabular}{p{0.155\linewidth}p{0.145\linewidth}p{0.145\linewidth}p{0.175\linewidth}p{0.135\linewidth}p{0.085\linewidth}}
\toprule
Harmonized variable & KNHANES (source/cycle) & NHANES (source/cycle) & Original item (abridged) & Transformation rule & Compar. \\
\midrule
Tinnitus presence & Otologic Q, 2009--2012 & Audiometry Q (AUQ), sel.\ cycles & ``Ringing/noise in ears past 12 mo?'' & Any vs.\ none (binary) & High \\
Bothersome tinnitus & Otologic Q, 2009--2012 & AUQ bother item, sel.\ cycles & Annoyance / sleep / concentration & Bothersome vs.\ not & Medium \\
Pure-tone thresholds & Audiometry exam, 2010--2012 & Audiometry exam, sel.\ cycles & Air-conduction dB HL per freq. & Better/worse-ear PTA; per-ear & High \\
Perceived stress & Health Q (stress item) & DPQ / stress-related item & ``How much stress in daily life?'' & Collapse to high vs.\ low & \textbf{Low} \\
Depressed mood & PHQ-9 / mood item & PHQ-9 (DPQ) & Standardized depression items & Screen-positive vs.\ negative & Medium \\
Occupation / employment & Socioecon.\ Q & Occupation (OCQ) & Job category, employed y/n & Common category crosswalk & Medium \\
Occupational noise & Sel.\ cycles & Noise exposure (AUQ) & Loud-noise exposure, duration & Duration/intensity + binary & Medium \\
Sleep duration & Health Q & Sleep (SLQ) & Average hours of sleep & Continuous hours & Medium \\
Smoking / alcohol & Health Q & SMQ / ALQ & Current status, quantity & Standard categories & High \\
Age, sex, education & Demographics & Demographics (DEMO) & --- & Direct / crosswalk & High \\
Survey weight, stratum, PSU & Design vars & Design vars & --- & Kept for design-consistent SEs & High \\
\bottomrule
\end{tabular}
\end{table}
 % Supplementary Table S3

\end{document}
